% ==============================================================================
% SECTION 3 — METHOD
% ==============================================================================
\section{Method}
\label{sec:method}

We apply the identical experimental protocol of Willis et al.\ \cite{Willis2025_llm_ipd}
to four new models, adding the first Google provider entry and enabling
cross-provider comparison. We summarise the key components below.

\subsection{Strategy Generation}

For each combination of model and prompt style, we generate 25 natural-language
\IPD strategies per attitude (Aggressive, Cooperative, Neutral), yielding 75
strategies per model--prompt pair. Three prompt styles are used (Table~\ref{tab:prompts}):
\textbf{Default} (direct elicitation with game-theoretic language),
\textbf{Refine} (Self-Refine \cite{Madaan2023_self_refine} applied to the default
output), and \textbf{Prose} (obfuscated scenario without game-theoretic terminology).

\begin{table}[h]
\caption{Prompt styles (following Willis et al.\ \cite{Willis2025_llm_ipd}).}
\label{tab:prompts}
\small
\begin{tabular}{@{}lp{5.5cm}@{}}
\toprule
\textbf{Style} & \textbf{Description} \\
\midrule
Default & Direct prompt with game-theoretic language; strategy generated in natural language. \\
Refine  & Default output refined via Self-Refine \cite{Madaan2023_self_refine}: the model critiques and rewrites its own strategy. \\
Prose   & Game-theoretic framing obfuscated as a real-world scenario (e.g., trade negotiation). Strategy generated, then translated to the \IPD context. \\
\bottomrule
\end{tabular}
\end{table}

Strategies are then converted to Python algorithms using the same \LLM that
generated them (one per provider), preserving consistency within each provider.
Strategies that cannot be executed without error are replaced by regeneration
until 25 valid strategies per attitude are obtained.

\subsection{Models}

We apply the protocol to the four models listed in Table~\ref{tab:models}.
These represent the direct successors of the models studied by Willis et al.\
(Claude lineage, OpenAI lineage) and a new provider (Google).

\begin{table}[h]
\caption{Models evaluated in this study.}
\label{tab:models}
\small
\begin{tabular}{@{}llll@{}}
\toprule
\textbf{Model} & \textbf{Provider} & \textbf{Predecessor} & \textbf{Released} \\
\midrule
Claude Sonnet 4.6    & Anthropic & Claude 3.5 Sonnet  & 2025 \\
Gemini 2.5 Flash     & Google    & (new provider)     & 2025 \\
Gemini 3.1 Pro       & Google    & (new provider)     & 2025 \\
GPT-5.4 Mini         & OpenAI    & ChatGPT-4o         & 2025 \\
\bottomrule
\end{tabular}
\end{table}

\subsection{IPD Tournament}

All 75 strategies compete in an all-play-all tournament using the Axelrod Python
library \cite{Knight2016_open}. Each match consists of 1,000 rounds of the
standard \IPD (payoff matrix: $R=3$, $S=0$, $T=5$, $P=1$). Noise conditions
introduce a 10\% probability of action-flip per player per round. Tournaments
are repeated 20 times.

\subsection{Attitude-Agents}

Following Willis et al.\ \cite{Willis2025_llm_ipd}, we define three attitude-agents,
each uniformly sampling from its corresponding strategy set for each match.
This captures populations of agents with distinct strategic dispositions rather
than fixed individual strategies.

\subsection{Moran Process}

We simulate Moran evolutionary processes with population size $n=12$ and 100
iterations per condition. Four population compositions are evaluated:
\begin{enumerate}
  \item \textbf{Balanced, clean} (4:4:4) --- equal priors;
  \item \textbf{Biased, clean} (8:2:2) --- aggressive majority;
  \item \textbf{Balanced, noise} (4:4:4 with noise) --- equal priors with action noise;
  \item \textbf{Biased, noise} (8:2:2 with noise) --- aggressive majority with action noise.
\end{enumerate}

Convergence is assessed by the proportion of runs reaching each monoculture
equilibrium (all-Aggressive, all-Cooperative, or all-Neutral).

\subsection{Derived Metrics}

We introduce the \emph{Index of Differential Capabilities} (\ICD) to summarise
the head-to-head payoff gap between aggressive and cooperative agents:
\begin{equation}
  \ICD = \frac{\bar{u}(A)}{\bar{u}(C)}, \quad
  \bar{u}(k) = \frac{1}{3}\sum_{j \in \{A,C,N\}} u(k,j),
  \label{eq:icd}
\end{equation}
where $u(k,j)$ is the normalised payoff of attitude $k$ against attitude $j$.
An \ICD of 1.0 implies equal capability; values below 1.0 indicate a cooperative
advantage. We also define the \emph{noise sensitivity} $\Dnoise$:
\begin{equation}
  \Dnoise = \PC^{\text{clean}} - \PC^{\text{noise}},
  \label{eq:dnoise}
\end{equation}
measuring the drop in cooperative equilibrium probability under action noise.
